% Options for packages loaded elsewhere
\PassOptionsToPackage{unicode}{hyperref}
\PassOptionsToPackage{hyphens}{url}
\PassOptionsToPackage{dvipsnames,svgnames,x11names}{xcolor}
\documentclass[
  10pt,
  fontset=fandol]{ctexart}
\usepackage{xcolor}
\usepackage[margin=16mm]{geometry}
\usepackage{amsmath,amssymb}
\setcounter{secnumdepth}{-\maxdimen} % remove section numbering
\usepackage{iftex}
\ifPDFTeX
  \usepackage[T1]{fontenc}
  \usepackage[utf8]{inputenc}
  \usepackage{textcomp} % provide euro and other symbols
\else % if luatex or xetex
  \usepackage{unicode-math} % this also loads fontspec
  \defaultfontfeatures{Scale=MatchLowercase}
  \defaultfontfeatures[\rmfamily]{Ligatures=TeX,Scale=1}
\fi
\usepackage{lmodern}
\ifPDFTeX\else
  % xetex/luatex font selection
\fi
% Use upquote if available, for straight quotes in verbatim environments
\IfFileExists{upquote.sty}{\usepackage{upquote}}{}
\IfFileExists{microtype.sty}{% use microtype if available
  \usepackage[]{microtype}
  \UseMicrotypeSet[protrusion]{basicmath} % disable protrusion for tt fonts
}{}
\makeatletter
\@ifundefined{KOMAClassName}{% if non-KOMA class
  \IfFileExists{parskip.sty}{%
    \usepackage{parskip}
  }{% else
    \setlength{\parindent}{0pt}
    \setlength{\parskip}{6pt plus 2pt minus 1pt}}
}{% if KOMA class
  \KOMAoptions{parskip=half}}
\makeatother
\setlength{\emergencystretch}{3em} % prevent overfull lines
\providecommand{\tightlist}{%
  \setlength{\itemsep}{0pt}\setlength{\parskip}{0pt}}
\usepackage{amsmath,amssymb,mathtools,needspace}
\xeCJKDeclareCharClass{CJK}{"2032,"0400->"04FF,"2160->"216B,"2460->"2473,"25A0,"25C7,"25CB,"25CF}
\IfFontExistsTF{Arial Unicode MS}{\xeCJKsetup{AutoFallBack=true}\setCJKfallbackfamilyfont{\CJKrmdefault}{Arial Unicode MS}}{}
\setcounter{secnumdepth}{0}
\setlength{\emergencystretch}{2em}
\allowdisplaybreaks
\usepackage{bookmark}
\IfFileExists{xurl.sty}{\usepackage{xurl}}{} % add URL line breaks if available
\urlstyle{same}
\hypersetup{
  pdftitle={误差分析},
  colorlinks=true,
  linkcolor={Maroon},
  filecolor={Maroon},
  citecolor={Blue},
  urlcolor={Blue},
  pdfcreator={LaTeX via pandoc}}

\title{误差分析}
\author{}
\date{}

\begin{document}
\maketitle

\subsection{7.6 误差分析}\label{ux8befux5deeux5206ux6790}

\subsubsection{7.6.1 矩阵的条件数}\label{ux77e9ux9635ux7684ux6761ux4ef6ux6570}

考虑线性方程组 \(Ax=b\)，其中设 \(A\) 为非奇异矩阵，\(x\) 为方程组的精确解。

由于 \(A\)（或 \(b\)）元素是测量得到的，或者是计算的结果，所以 \(A\)（或 \(b\)）常带有某些观测误差或者包含有舍入误差。因此，在处理实际问题 \((A+\delta A)x=b+\delta b\) 时，通常要考虑 \(A\)（或 \(b\)）的微小误差对解的影响，即考虑估计 \(x-y\)，其中 \(y\) 是 \((A+\delta A)y=b+\delta b\) 的解。

首先考察一个例子。

\textbf{例 7.8} 设有方程组

\[
\begin{pmatrix}1&1\\1&1.0001\end{pmatrix}
\begin{pmatrix}x_1\\x_2\end{pmatrix}
=\begin{pmatrix}2\\2\end{pmatrix}
\]

或

\[
Ax=b,
\tag{7.6.1}
\]

它的精确解为

\[
x=(2,0)^{\mathrm T}.
\]

现在考虑常数项的微小变化对方程组解的影响，即考察方程组

\[
\begin{pmatrix}1&1\\1&1.000\,1\end{pmatrix}
\begin{pmatrix}y_1\\y_2\end{pmatrix}
=\begin{pmatrix}2\\2.000\,1\end{pmatrix}
\]

或

\[
A(x+\delta x)=b+\delta b,
\tag{7.6.2}
\]

其中 \(\delta b=(0,0.000\,1)^{\mathrm T},y=x+\delta x\)，\(x\) 为式（7.6.1）的解。显然方程组（7.6.2）的解为

\[
x+\delta x=(1,1)^{\mathrm T}.
\]

可以看到，式（7.6.1）的常数项 \(b\) 的第二个分量只有 \(\frac1{10\,000}\) 的微小变化，方程组的解却变化很大。这样的方程组称为病态方程组。

\textbf{定义 7.7} 如果矩阵 \(A\) 或常数项 \(b\) 的微小变化，引起方程组 \(Ax=b\) 解的巨大变化，则称此方程组为\textbf{病态方程组}，矩阵 \(A\) 称为\textbf{病态矩阵}（相对于方程组而言），否则称方程组为\textbf{良态方程组}，\(A\) 称为\textbf{良态矩阵}。

应该注意，矩阵的病态性质是矩阵本身的特性，下面希望找出刻画矩阵病态性质

（正文续下页。）

\begin{quote}
校注（agent 补充，原书疑误）：本页定理 7.18 证明原文为 \(Bx_0=x_0\)，以及 \((I\pm B)^{-1}=I\mp B(I-B)^{-1}\)，均按印刷内容保留。若同时处理 \(I+B\) 和 \(I-B\) 两种情形，前式应为 \(Bx_0=\mp x_0\)，后式右端的逆矩阵应为 \((I\pm B)^{-1}\)；原文写法只与减号分支一致。前式符号不影响随后所取范数等于 1 的结论。\href{../assets/ch07b/205-inverse-signs.png}{原式局部图}。
\end{quote}

\href{../../source-images/pdf-206.jpeg}{原页图像}

（续 7.6.1 矩阵的条件数。）

的量。设有方程组

\[
Ax=b,
\tag{7.6.3}
\]

其中 \(A\) 为非奇异阵，\(x\) 为式（7.6.3）的精确解。下面研究当方程组的系数矩阵 \(A\)（或 \(b\)）有微小误差（扰动）时对解的影响。

现设 \(A\) 是精确的，\(b\) 有误差 \(\delta b\)，解为 \(x+\delta x\)，则

\[
A(x+\delta x)=b+\delta b,\quad\delta x=A^{-1}\delta b,\quad\|\delta x\|\leqslant\|A^{-1}\|\|\delta b\|.
\tag{7.6.4}
\]

由式（7.6.3），有

\[
\|b\|\leqslant\|A\|\|x\|,
\]

即

\[
\frac1{\|x\|}\leqslant\frac{\|A\|}{\|b\|}\quad\text{（设 }b\neq0\text{）},
\tag{7.6.5}
\]

于是由式（7.6.4）式及式（7.6.5），有如下定理。

\textbf{定理 7.19} 设 \(A\) 是非奇异阵，\(Ax=b\neq0\)，且 \(A(x+\delta x)=b+\delta b\)，则

\[
\frac{\|\delta x\|}{\|x\|}\leqslant\|A^{-1}\|\|A\|\frac{\|\delta b\|}{\|b\|}.
\]

上式给出了解的相对误差的上界，常数项 \(b\) 的相对误差在解中可能放大 \(\|A^{-1}\|\|A\|\) 倍。

现设 \(b\) 是精确的，\(A\) 有微小误差（扰动）\(\delta A\)，解为 \(x+\delta x\)，则

\[
(A+\delta A)(x+\delta x)=b,
\]

\[
(A+\delta A)\delta x=-(\delta A)x.
\tag{7.6.6}
\]

如果 \(\delta A\) 不受限制的话，\(A+\delta A\) 可能奇异，而

\[
(A+\delta A)=A(I+A^{-1}\delta A).
\]

由定理 7.18 知，当 \(\|A^{-1}\delta A\|<1\) 时，\((I+A^{-1}\delta A)^{-1}\) 存在，由式（7.6.6），有

\[
\delta x=-(I+A^{-1}\delta A)^{-1}A^{-1}(\delta A)x,\quad
\|\delta x\|\leqslant\frac{\|A^{-1}\|\|\delta A\|\|x\|}{1-\|A^{-1}(\delta A)\|}.
\]

设 \(\|A^{-1}\|\|\delta A\|<1\)，即得

\[
\frac{\|\delta x\|}{\|x\|}
\leqslant
\frac{\|A^{-1}\|\|A\|\dfrac{\|\delta A\|}{\|A\|}}
{1-\|A^{-1}\|\|A\|\dfrac{\|\delta A\|}{\|A\|}}.
\tag{7.6.7}
\]

\textbf{定理 7.20} 设 \(A\) 为非奇异矩阵，\(Ax=b\neq0\)，且 \((A+\delta A)(x+\delta x)=b\)，如果 \(\|A^{-1}\|\|\delta A\|<1\)，则式（7.6.7）成立。

如果 \(\delta A\) 充分小，且在 \(\|A^{-1}\|\|\delta A\|<1\) 条件下，那么式（7.6.7）说明矩阵 \(A\) 的相对误差 \(\frac{\|\delta A\|}{\|A\|}\) 在解中可能放大 \(\|A^{-1}\|\|A\|\) 倍。

总之，量 \(\|A^{-1}\|\|A\|\) 愈小，由 \(A\)（或 \(b\)）的相对误差引起的解的相对误差就愈小；量 \(\|A^{-1}\|\|A\|\) 愈大，该解的相对误差就可能愈大。所以量 \(\|A^{-1}\|\|A\|\) 实际上刻画了解对原始数据变化的灵敏度，即刻画了方程组的病态程度，于是引进下述定

（定义续下页。）

\href{../../source-images/pdf-207.jpeg}{原页图像}

（续 7.6.1 矩阵的条件数。）

义。

\textbf{定义 7.8} 设 \(A\) 为非奇异矩阵，称数 \(\operatorname{cond}(A)_v=\|A^{-1}\|_v\|A\|_v\)（\(v=1,2\) 或 \(\infty\)）为矩阵 \(A\) 的\textbf{条件数}。

由此看出矩阵的条件数与范数有关。

矩阵的条件数是一个十分重要的概念。由上面讨论知，当 \(A\) 的条件数相对地大时，即 \(\operatorname{cond}(A)\gg1\) 时，式（7.6.3）是病态的（即 \(A\) 是病态矩阵，或者说 \(A\) 是坏条件的）；当 \(A\) 的条件数相对地小时，式（7.6.3）是良态的（或者说 \(A\) 是好条件的）。\(A\) 的条件数愈大，方程组的病态程度愈严重，也就愈难得到方程组的比较准确的解。

通常使用的条件数有

（1）\(\operatorname{cond}(A)_\infty=\|A^{-1}\|_\infty\|A\|_\infty\)；

（2）\(A\) 的\textbf{谱条件数}

\[
\operatorname{cond}(A)_2=\|A\|_2\|A^{-1}\|_2
=\sqrt{\frac{\lambda_{\max}(A^{\mathrm T}A)}{\lambda_{\min}(A^{\mathrm T}A)}}.
\]

当 \(A\) 为对称阵时，\(\operatorname{cond}(A)_2=\frac{|\lambda_1|}{|\lambda_n|}\)，其中 \(\lambda_1,\lambda_n\) 为 \(A\) 的绝对值最大和绝对值最小的特征值。

条件数有如下性质：

（1）对任何非奇异矩阵 \(A\)，都有 \(\operatorname{cond}(A)_v\geqslant1\)。事实上，

\[
\operatorname{cond}(A)_v=\|A^{-1}\|_v\|A\|_v\geqslant\|A^{-1}A\|_v=1.
\]

（2）设 \(A\) 为非奇异矩阵，\(c\) 为不等于零的常数，则 \(\operatorname{cond}(cA)_v=\operatorname{cond}(A)_v\)。

（3）如果 \(A\) 为正交矩阵，则 \(\operatorname{cond}(A)_2=1\)；如果 \(A\) 为非奇异矩阵，\(R\) 为正交矩阵，则

\[
\operatorname{cond}(RA)_2=\operatorname{cond}(AR)_2=\operatorname{cond}(A)_2.
\]

\textbf{例 7.9} 已知 Hilbert 矩阵

\[
H_n=\begin{pmatrix}
1&\frac12&\cdots&\frac1n\\[4pt]
\frac12&\frac13&\cdots&\frac1{n+1}\\
\vdots&\vdots&&\vdots\\
\frac1n&\frac1{1+n}&\cdots&\frac1{2n-1}
\end{pmatrix},
\]

计算 \(H_3\) 的条件数。

\textbf{解}

\[
H_3=\begin{pmatrix}
1&\frac12&\frac13\\[4pt]
\frac12&\frac13&\frac14\\[4pt]
\frac13&\frac14&\frac15
\end{pmatrix},\quad
H_3^{-1}=\begin{pmatrix}
9&-36&30\\
-36&192&-180\\
30&-180&180
\end{pmatrix}.
\]

（例 7.9 解续下页。）

\href{../../source-images/pdf-208.jpeg}{原页图像}

（续 7.6.1 例 7.9。）

（1）计算 \(H_3\) 条件数 \(\operatorname{cond}(H_3)_\infty\)。

\[
\|H_3\|_\infty=11/6,\quad\|H_3^{-1}\|_\infty=408,
\]

所以 \(\operatorname{cond}(H_3)_\infty=748\)。同样可计算 \(\operatorname{cond}(H_6)=2.9\times10^6\)。一般当 \(n\) 愈大时，\(H_n\) 的病态愈严重。

（2）考虑方程组

\[
H_3x=(11/6,13/12,47/60)^{\mathrm T}=b.
\]

设 \(H_3\) 及 \(b\) 有微小误差（取三位有效数字），有

\[
\begin{pmatrix}
1.00&0.500&0.333\\
0.500&0.333&0.250\\
0.333&0.250&0.200
\end{pmatrix}
\begin{pmatrix}x_1+\delta x_1\\x_2+\delta x_2\\x_3+\delta x_3\end{pmatrix}
=\begin{pmatrix}1.83\\1.08\\0.783\end{pmatrix},
\tag{7.6.8}
\]

简记作 \((H_3+\delta H_3)(x+\delta x)=b+\delta b\)。方程组 \(H_3x=b\) 与式（7.6.8）的精确解为

\[
x=(1,1,1)^{\mathrm T},\quad
x+\delta x=(1.089\,512\,538,\ 0.487\,967\,062,\ 1.491\,002\,798)^{\mathrm T},
\]

于是

\[
\delta x=(0.089\,5,\ -0.512\,0,\ 0.491\,0)^{\mathrm T},\quad
\frac{\|\delta H_3\|_\infty}{\|H_3\|_\infty}\approx0.18\times10^{-3}<0.02\%,
\]

\[
\frac{\|\delta b\|_\infty}{\|b\|_\infty}\approx0.182\%,\quad
\frac{\|\delta x\|_\infty}{\|x\|_\infty}\approx51.2\%.
\]

这就是说，\(H_3\) 与 \(b\) 相对误差不超过 \(0.3\%\)，而引起解的相对误差超过 \(50\%\)。

由上面讨论可知，要判别一个矩阵是否病态，需要计算条件数 \(\operatorname{cond}(A)=\|A^{-1}\|\|A\|\)，而计算 \(A^{-1}\) 是比较费劲的。那么在实际计算中如何发现病态情况呢？

（1）如果在 \(A\) 的三角约化时（尤其是用主元素消去法解方程组（7.6.3）时）出现小主元，那么对大多数矩阵来说，\(A\) 是病态矩阵。例如用选主元的直接三角分解法解方程组（7.6.8）（用双精度累加计算 \(\sum_i a_ib_i\)，结果舍入为三位浮点数），则有

\[
I_{23}(H_3+\delta H_3)
=\begin{pmatrix}
1&&\\
0.333&1&\\
0.500&0.994&1
\end{pmatrix}
\begin{pmatrix}
1&0.500\,0&0.333\,0\\
&0.083\,5&0.089\,1\\
&&\boxed{-0.005\,07}
\end{pmatrix}
=LU.
\]

（2）如果 \(A\) 的最大特征值和最小特征值之比（按绝对值）是大的，则 \(A\) 是病态的，即当 \(|\lambda_1|/|\lambda_n|\) 是大的，其中 \(A\) 特征值次序为

\[
|\lambda_1|\geqslant|\lambda_2|\geqslant\cdots\geqslant|\lambda_n|>0.
\]

显然

\[
|\lambda_1|\leqslant\|A\|,\quad\frac1{|\lambda_n|}\leqslant\|A^{-1}\|,
\]

因此

\[
\operatorname{cond}(A)\geqslant\frac{|\lambda_1|}{|\lambda_n|}\gg1.
\]

（3）如果系数矩阵的行列式值相对来说很小，或系数矩阵某些行近似线性相关，则 \(A\) 可能是病态的。

（条目（4）续下页。）

\begin{quote}
校注（agent 补充，原书疑误）：原文印作 \(\operatorname{cond}(H_6)=2.9\times10^6\)，这里保留。沿用本例的 \(\infty\)-范数，对 \(6\times6\) Hilbert 矩阵作有理数精确运算得到 \(\|H_6\|_\infty\|H_6^{-1}\|_\infty=29\,070\,279\approx2.9\times10^7\)，原书指数疑应为 7。\href{../assets/ch07b/208-h6-cond.png}{原文局部图}。
\end{quote}

\begin{quote}
校注（agent 补充，ch07b-E014）：本页称为``精确解''的小数向量按原书照录。对式（7.6.8）所列有限小数作精确有理数求解，结果为 \((153\,971,\allowbreak 68\,960,\allowbreak 210\,710)^{\mathrm T}/141\,321\approx(1.0895125282,\allowbreak 0.4879671103,\allowbreak 1.4910027526)^{\mathrm T}\)，与所印三个分量的末几位不符。因此原书这些小数不能视作式（7.6.8）的精确解；该差异不影响后文约 \(51.2\%\) 的结论。
\end{quote}

\href{../../source-images/pdf-209.jpeg}{原页图像}

（续 7.6.1 矩阵的条件数。）

（4）如果系数矩阵 \(A\) 元素间数量级相差很大，并且无一定规则，则 \(A\) 可能是病态的。

病态问题通常不能用选主元素的消去法来解决。对于此类问题一般采用高精度的算术运算（采用双倍字长进行运算）或者预处理方法。即将求解 \(Ax=b\) 的问题转化为求解一等价方程组

\[
\begin{cases}
PAQy=Pb,\\
y=Q^{-1}x.
\end{cases}
\]

通过选择非奇异矩阵 \(P,Q\)，使 \(\operatorname{cond}(PAQ)<\operatorname{cond}(A)\)。一般选择 \(P,Q\) 为对角阵或者三角阵。

当矩阵 \(A\) 的元素大小不均时，在 \(A\) 的行（或列）中引进适当的比例因子（使矩阵 \(A\) 的所有行或列按 \(\infty\)-范数大体上有相同的长度，使 \(A\) 的系数均衡），对 \(A\) 的条件数是有影响的。但这种方法不能保证 \(A\) 的条件数一定得到改善。

\textbf{例 7.10} 设 \(A=\begin{pmatrix}1&10^4\\1&1\end{pmatrix}\)，且

\[
\begin{pmatrix}1&10^4\\1&1\end{pmatrix}
\begin{pmatrix}x_1\\x_2\end{pmatrix}
=\begin{pmatrix}10^4\\2\end{pmatrix},
\tag{7.6.9}
\]

计算 \(\operatorname{cond}(A)_\infty\)。

\textbf{解} 由

\[
A=\begin{pmatrix}1&10^4\\1&1\end{pmatrix},
\]

有

\[
A^{-1}=\frac1{10^4-1}\begin{pmatrix}-1&10^4\\1&-1\end{pmatrix},\quad
\operatorname{cond}(A)_\infty=\frac{(1+10^4)^2}{10^4-1}\approx10^4.
\]

现在 \(A\) 的第 1 行引进比例因子。如用 \(s_1=\max_{1\leqslant i\leqslant2}|a_{1i}|=10^4\) 除第一个方程式，得 \(A'x=b'\)，即

\[
\begin{pmatrix}10^{-4}&1\\1&1\end{pmatrix}
\begin{pmatrix}x_1\\x_2\end{pmatrix}
=\begin{pmatrix}1\\2\end{pmatrix},
\tag{7.6.10}
\]

\[
(A')^{-1}=\frac1{1-10^{-4}}\begin{pmatrix}-1&1\\1&-10^{-4}\end{pmatrix},\quad
\operatorname{cond}(A')_\infty=\frac1{1-10^{-4}}\approx4.
\]

当用列主元素消去法解式（7.6.9）时（计算到三位数字），得

\[
(A,b)\longrightarrow
\left(\begin{array}{cc|c}
1&10^4&10^4\\
0&-10^4&-10^4
\end{array}\right),
\]

于是得到很坏的结果：\(x_1=0,x_2=1\)。

现用列主元素消去法解式（7.6.10），得

（例 7.10 解续下页。）

\begin{quote}
校注（agent 补充，原书疑误）：\(\operatorname{cond}(A')_\infty\) 的分式分子原图为 1，而后面印作 \(\approx4\)，这里保留。由本页两个矩阵，\(\|A'\|_\infty=2\)，\(\|(A')^{-1}\|_\infty=2/(1-10^{-4})\)，故条件数应为 \(4/(1-10^{-4})=40\,000/9\,999\approx4.0004\)。\href{../assets/ch07b/209-scaled-cond.png}{原式局部图}。
\end{quote}

\href{../../source-images/pdf-210.jpeg}{原页图像}

（续 7.6.1 例 7.10。）

\[
(A',b')\longrightarrow
\left(\begin{array}{cc|c}1&1&2\\10^{-4}&1&1\end{array}\right)
\longrightarrow
\left(\begin{array}{cc|c}1&1&2\\0&1&1\end{array}\right),
\]

从而得到较好的结果：\(x_1=1,x_2=1\)。

设 \(\overline x\) 为方程组 \(Ax=b\) 的近似解，于是可计算 \(\overline x\) 的剩余向量 \(r=b-A\overline x\)，当 \(r\) 很小时，\(\overline x\) 是否为 \(Ax=b\) 一个较好的近似解呢？下述定理给出了解答。

\textbf{定理 7.21（事后误差估计）}

\(1^\circ\) 设 \(A\) 为非奇异矩阵，\(x\) 是精确解，\(Ax=b\neq0\)；

\(2^\circ\) 设 \(\overline x\) 是方程组的近似解，\(r=b-A\overline x\)，则

\[
\frac{\|x-\overline x\|}{\|x\|}\leqslant\operatorname{cond}(A)\cdot\frac{\|r\|}{\|b\|}.
\tag{7.6.11}
\]

\textbf{证明} 由 \(x-\overline x=A^{-1}r\)，得

\[
\|x-\overline x\|\leqslant\|A^{-1}\|\|r\|,
\tag{7.6.12}
\]

又有

\[
\|b\|=\|Ax\|\leqslant\|A\|\|x\|,
\]

\[
\frac1{\|x\|}\leqslant\frac{\|A\|}{\|b\|},
\tag{7.6.13}
\]

由式（7.6.12）及式（7.6.13）即得到式（7.6.11）。

\subsubsection{7.6.2 舍入误差}\label{ux820dux5165ux8befux5dee}

在复杂的计算中，由浮点运算而引进的舍入误差可能积累而影响答案，因此对任何算法都需要进行舍入误差分析，看其是否过度影响所得的结果。

下面叙述采用选主元素 Gauss 消去法解式（7.6.3）的计算过程中舍入误差对解的影响。在 Gauss 消去法的舍入误差的分析和研究方面，von Neumann（在 1947 年）、Givens（在 1954 年）和 Wilkinson（在 1961 年、1963 年）等人都作出了自己的贡献。

设 \(\overline x\) 为用选主元素 Gauss 消去法解式（7.6.3）的计算解，\(x\) 为式（7.6.3）的精确解。若要直接计算每一步舍入误差对解的影响来获得界的估计 \(\|x-\overline x\|\)，那将是非常困难的。Wilkinson 等人提出了``向后误差分析方法''，其基本思想是把计算过程中舍入误差对解的影响归结为原始数据变化对解的影响，即计算解 \(\overline x\) 是下述扰动方程组的精确解 \((A+\delta A)x=b\)，其中 \(\delta A\) 为某个``小''矩阵。

下面给出一个定理来说明这个结果。

\textbf{定理 7.22（选主元素 Gauss 消去法误差分析）} 如果

\(1^\circ\) 设 \(A\) 为 \(n\) 阶非奇异矩阵；

\(2^\circ\) 用列主元素消去法（或完全主元素消去法）解式（7.6.3）；

（定理条件与结论续下页。）

\begin{quote}
校注（agent 补充，原书疑误）：向后误差分析段中印刷方程为 \((A+\delta A)x=b\)，而紧邻文字明确指计算解 \(\overline x\) 是其精确解。此处保留原式；按本段符号约定，应写 \((A+\delta A)\overline x=b\)。\href{../assets/ch07b/210-backward-error.png}{原文局部图}。
\end{quote}

\href{../../source-images/pdf-211.jpeg}{原页图像}

（续 7.6.2 定理 7.22。）

\(3^\circ\) 记 \(a_k=\max_{1\leqslant i,j\leqslant n}|a_{ij}^{(k)}|\)，\(a=\max_{1\leqslant i,j\leqslant n}|a_{ij}|\)，\(r=\max_{1\leqslant k\leqslant n}a_k/a\)------元素的增长因子，\(A^{(k)}\equiv(a_{ij}^{(k)})_n\)；

\(4^\circ\) \(t\) 为计算机字长（指尾数部分），矩阵的阶数 \(n\) 满足 \(n\cdot2^{-t}\leqslant0.01\)，

则（1）用选主元素 Gauss 消去法计算的三角阵 \(L,U\) 满足 \(LU=A+E\)，其中

\[
|(E)_{ij}|\leqslant2(n-1)ra2^{-t};
\]

（2）用选主元素 Gauss 消去法得到的计算解 \(\overline x\) 精确满足 \((A+\delta A)x=b\)，其中

\[
\|\delta A\|_\infty\leqslant1.01(n^3+3n^2)ar2^{-t}\leqslant1.01(n^3+3n^2)r\|A\|_\infty2^{-t};
\]

（3）计算解精度估计（\(x\) 为 \(Ax=b\) 精确解）

\[
\frac{\|x-\overline x\|_\infty}{\|x\|_\infty}
\quad
\frac{\operatorname{cond}(A)_\infty}
{1-\operatorname{cond}(A)_\infty\dfrac{\|\delta A\|_\infty}{\|A\|_\infty}}
\times1.01(n^3+3n^2)r2^{-t}.
\tag{7.6.14}
\]

上述计算解 \(\overline x\) 精度估计式说明，\(\overline x\) 的相对误差限依赖于 \(\operatorname{cond}(A)_\infty\)、元素的增长因子、方程组阶数、计算机字长等。

\begin{center}\rule{0.5\linewidth}{0.5pt}\end{center}

\subsection{相关原页校注（agent 补充）}\label{ux76f8ux5173ux539fux9875ux6821ux6ce8agent-ux8865ux5145}

以下校注来自本节覆盖的原页，可能也涉及同页相邻小节；原文照录与校正说明分开保留。

\subsubsection{原书 PDF 页 211 的校注}\label{ux539fux4e66-pdf-ux9875-211-ux7684ux6821ux6ce8}

\href{../pages/pdf-211.md}{完整原页转录} · \href{../../source-images/pdf-211.jpeg}{原页图像}

校注记录：ch07b-E011。

\begin{quote}
校注（agent 补充，原书疑误）：定理 7.22（2）印作``计算解 \(\overline x\) 精确满足 \((A+\delta A)x=b\)''，此处按原文保留；与上一页同样，方程中的 \(x\) 应按符号约定写为 \(\overline x\)。\href{../assets/ch07b/211-computed-solution.png}{原式局部图}。

校注（agent 补充，原书疑误）：式（7.6.14）原图在左侧误差比与右侧分式之间未印关系符号，此处保持两式并列；按精度估计用途应为 \(\leqslant\)。另外，由式（7.6.7）得到此界还须 \(\operatorname{cond}(A)_\infty\|\delta A\|_\infty/\|A\|_\infty<1\)，保证分母为正；这一适用条件未在本定理列出。\href{../assets/ch07b/211-error-estimate.png}{原式局部图}。
\end{quote}

\subsection{本节覆盖原页的校注记录}\label{ux672cux8282ux8986ux76d6ux539fux9875ux7684ux6821ux6ce8ux8bb0ux5f55}

下列记录按原页关联，含同页相邻小节的校注；计算时同时读取上文校注和完整原页。

\begin{itemize}
\tightlist
\item
  \texttt{ch07b-E007}：原书 PDF 页 205
\item
  \texttt{ch07b-E008}：原书 PDF 页 208
\item
  \texttt{ch07b-E009}：原书 PDF 页 209
\item
  \texttt{ch07b-E010}：原书 PDF 页 210
\item
  \texttt{ch07b-E011}：原书 PDF 页 211
\item
  \texttt{ch07b-E014}：原书 PDF 页 208
\end{itemize}

\end{document}
